\begin{table}[htbp]
\centering
\caption{Number of repeats certifying each lower bound, with two ensembles of \(N\) independent endpoints per repeat. Each row contains 200 repeats. A count of one is the trivial lower bound.}
\label{tab:counts}
\begin{tabular}{lr rrrr}
\toprule
&&\multicolumn{4}{c}{Certified lower bound}\\
\cmidrule(l){3-6}
Data source & \(N\) & 1 & 2 & 3 & 4\\
\midrule
Diffusion & 1,000   & 200 & 0   & 0   & 0\\
           & 10,000  & 0   & 0   & 114 & 86\\
           & 100,000 & 0   & 0   & 0   & 200\\
\addlinespace
Three-state hub & 1,000   & 200 & 0   & 0   & 0\\
                & 10,000  & 1   & 167 & 32  & 0\\
                & 100,000 & 0   & 0   & 200 & 0\\
\bottomrule
\end{tabular}
\end{table}

Table~\ref{tab:counts} uses the fixed isotropic \(n=18,w=1/3,t=1/128\) design and \(\alpha=0.05\). We draw endpoint counts directly from the multinomial and binomial laws above. The three-state control takes the same \(a_i\) as the diffusion and attains \(\E_3\) by Theorem~\ref{thm:finite}, placing it on the exclusion boundary. The diffusion's population lower bound of four was recovered in 86 of 200 repeats at 10,000 trials per ensemble and in all 200 at 100,000. None of the 600 control repeats falsely excluded three. The probability guarantee follows from the joint limits; these simulated counts describe power at the stated design points.
